\chapter{Conclusões, limitações e próximos passos}\label{cap:conclusao}
\rubricaComunicacao
A segunda etapa do projeto Higrolar cumpriu o propósito de integrar o protótipo físico da AP1 à infraestrutura de comunicação e construir um dataset real e rastreável. A arquitetura baseada em ESP32, DHT11, OLED, Wi-Fi e EMQX Cloud ofereceu publicações MQTT com QoS 1, status retido, testamento \texttt{offline}, configuração remota do intervalo e um assinante Paho que persiste mensagens em disco. A segregação em três usuários MQTT documentou responsabilidades diferentes de dispositivo, monitor e gravador, embora o modo de ACL utilizado no EMQX tenha sido permissivo por padrão.
A campanha definitiva de 25h29min27s totalizou 3.096 mensagens, sendo 3.082 telemetrias e 3.080 telemetrias únicas, com duas duplicações lógicas, uma medição sem timestamp NTP e nenhum salto interno detectável entre sequências de um mesmo boot. A temperatura registrada variou de 26,9 a 36,9~\si{\celsius} e a umidade relativa entre 34\% e 73\%. O primeiro acionamento do umidificador coincidiu com aumento da média de UR de 34,15\% para 52,87\% em janelas adjacentes de 30~min. O ensaio posterior com ar-condicionado e umidificador simultâneos (12h59--14h10) registrou médias de 30,14~\si{\celsius} e 59,28\% de UR, enquanto um segundo trecho simultâneo ocorreu das 15h02 às 15h19. Essas observações não estabelecem causalidade porque os fatores ambientais não foram rigidamente controlados.
\section{Limitações observadas}
O DHT11 não é sensor metrológico de alta precisão e não é adequado para quantificar umidades extremamente baixas; o ponto físico não representa necessariamente todo o cômodo. A fila de 120 registros em RAM não resiste à falta de energia, embora o \texttt{boot\_id} sobreviva em NVS. O QoS 1 admite repetição, e 100\% de completude entre sequências numeradas não significa coleta perfeita entre todos os períodos operacionais. Durante a campanha, uma transição de estado em 24/10 às 11h47 não teve causa documentada. O atraso de chegada é influenciado por clocks e backlog. O custo real do rádio não foi medido; a autonomia permanece estimada com base na AP1, sem atualização instrumental da corrente de Wi-Fi. O gravador valida o certificado TLS no Windows, mas o firmware da campanha usa \texttt{setInsecure()} sem CA embarcada; a ACL do EMQX estava em modo blacklist.
\section{Perspectivas para o Módulo 3}
Como consequência prática, a equipe dispõe de um fluxo funcional do sensor ao CSV, de registros de estado e de eventos do mundo real com horário. No ciclo seguinte poderão ser criados rótulos exploratórios de intervenção, testadas regras de limiares e calculadas características temporais, preservando a diferença entre medição, evento relatado e inferência. Recomendam-se campanhas repetidas em dias distintos, identificação da posição do nó e dos equipamentos, endurecimento da ACL para negação por padrão, configuração de CA no ESP32 e medição direta de consumo elétrico. O firmware, gravador, analisador, diário, CSV bruto e fontes LaTeX integram o pacote entregue para reprodutibilidade.
\section{Divisão do trabalho}
Os papéis do \moduloDoCiclo\ rotacionaram em relação à AP1: ninguém repetiu o papel do ciclo anterior.
\begin{description}[style=unboxed, leftmargin=0.6cm]
  \item[Raimundo Nonato da Silva Sá Sobrinho (\emph{firmware} e comunicação; na AP1, integração).] Firmware de comunicação: conexão Wi-Fi e MQTT, reconexão, fila local, NTP e testamento.
  \item[Francisco de Assis Machado dos Santos (\emph{firmware} e comunicação; na AP1, documentação).] Contrato de dados e de tópicos, configuração remota do período e contas do EMQX.
  \item[Lemuel Cavalcante Lopes (integração; na AP1, \emph{firmware}).] Caminho do dado do sensor ao CSV, matriz de decisão de conectividade e rastreabilidade da Seção~\ref{sec:rastreabilidade}.
  \item[Zacarias do Nascimento Seriano Junnior (dados; na AP1, \emph{hardware}).] Gravador, condução da campanha de coleta, diário e registro de eventos, e execução do \texttt{qualidade.py}.
  \item[Thales Rodrigues Martins (documentação; na AP1, integração).] Redação e revisão deste relatório, descrição do conjunto de dados e declaração do Apêndice~\ref{ap:ia}.
\end{description}
\bigskip
\imprimirTabelaIntegrantes
